\setlength{\baselineskip}{20pt}
\begin{abstract}

扩散模型作为一类强大的生成式框架，通过逐步添加高斯噪声并训练去噪网络实现反向预测，已在图像生成、音频合成等领域展现出卓越的性能。然而，扩散模型在判别式任务中的潜力尚未得到充分挖掘。表征学习作为判别式视觉任务的核心环节，旨在学习具有高判别性的数据特征表示，以服务于图像分类、目标检测和图像分割等下游任务。目前，已有少数研究尝试将扩散机制引入特定的判别式场景，但这些方法通常针对单一任务进行定制化设计，缺乏通用性。此外，在推理阶段引入递归式的去噪操作会带来显著的额外计算开销。因此，如何将扩散模型的去噪能力以通用、高效的方式融入判别式表征学习，在不增加推理成本的前提下提升特征的判别性，成为一个亟待解决的关键问题。

本文围绕基于扩散模型的图像判别任务特征增强算法展开系统研究，研究内容和主要贡献如下：

（1）面向无额外推理成本的特征增强算法研究。本文创新性地将骨干网络中每个嵌入层重新解释为去噪层，将级联的$N$个嵌入层类比为扩散模型中$T$步递归去噪过程，在统一框架下实现特征提取与特征去噪的联合建模。在此基础上，通过严格的数学推导提出参数融合策略，证明去噪层参数可等价地合并至原始嵌入层中，使推理阶段的特征去噪无需任何额外计算开销。该方法无需标签信息即可提升特征质量，同时在有标签时亦可作为互补机制进一步增强性能。在行人重识别、图像分类、目标检测和图像分割等多项视觉任务上的实验结果表明，该方法能够稳定且显著地提升各类基线模型的性能。

（2）基于模型蒸馏的特征增强算法研究。针对轻量化去噪模块表达能力受限以及有限去噪深度下效果不足的问题，本文进一步提出两项增强策略。一方面，引入教师—学生蒸馏框架，利用高容量的教师网络为轻量化的学生去噪模块提供更精确的噪声分布估计指导，有效校准去噪方向并降低预测偏差，且该额外开销仅存在于训练阶段。另一方面，提出多步去噪融合策略，包含多轮单层去噪与基于渐进蒸馏的多步融合两种互补机制，通过数学推导将多步迭代过程折叠为等效的单步线性运算，在有限层数下显著增强去噪能力。两项策略分别从噪声预测精度与等效去噪幅度两个维度互补增强，共同实现了更优的特征判别性。大量实验验证了所提方法在多种骨干架构上的有效性与泛化能力，且始终保持推理阶段零额外计算开销的核心优势。

\par\noindent\keywords{扩散模型；表征学习；特征去噪；参数融合；知识蒸馏}
\end{abstract}